\begin{frame}{사전학습 vs 파인튜닝}
  \small
  \begin{tabular}{@{}p{1.5cm}p{5.0cm}p{5.0cm}@{}}
    \toprule
     & 사전학습(Pretraining) & 파인튜닝(Fine-tuning) \\
    \midrule
    데이터 & 인터넷 규모의 방대한 텍스트 & 특정 목적에 맞는 비교적 적은 데이터 \\
    목표 & 다음 토큰 예측(비지도) & 특정 작업(대화, 지시 따르기 등)에 맞춘 지도학습 \\
    결과 & 폭넓지만 다듬어지지 않은 능력 & 목적에 맞게 다듬어진 행동 \\
    \bottomrule
  \end{tabular}
  \vskip1em
  \begin{itemize}
    \item 사전학습: 계산 비용 \textbf{막대} --- 수천 GPU,
      몇 주$\sim$몇 달. 파인튜닝: 상대적으로 적은 자원.
    \item 핵심: ``\textbf{기초를 다시 배우지 않고, 이미 아는
      것을 특정 방향으로 다듬는다}'' --- Week 10.3의
      \kb{전이학습}(사전학습 CNN 앞쪽 층 고정, 마지막 층만
      재학습)과 \textbf{정확히 같은 논리}. 이미지 대신
      텍스트, 합성곱 필터 대신 Transformer 가중치.
  \end{itemize}
\end{frame}
